
In order to analyze the string usage in the scraped repositories, we first create and populate the tables featured in the
repository to then execute the select queries.
We decided to use DuckDB~\cite{raasveldt_duckdb_2019} as the database system, as it is fast
and easy to set up locally.

In total, we retrieved 36,991,829 queries from 27,389 repositories,
with 752,655 \texttt{SELECT}. By parsing the \texttt{CREATE} SQL statements, we could extract
701,382 distinct tables with 4,794,784 columns. In order to execute the queries
later in DuckDB, we unified the system specific column types into abstract types like \texttt{INTEGER}, \texttt{FLOAT}, 
\texttt{TEXT} or \texttt{DATETIME} to then later map them to the DuckDB types. 

For each repository, we then created a DuckDB database and created the tables using the parsed \texttt{CREATE} statements.
We then executed all the \texttt{INSERT} statements to populate the tables with data. Due to errors in the extracted 
SQL and differences in the features, data types and SQL dialects across systems, we could only execute 86.33\% 
of the \texttt{CREATE TABLE}, 83.83\% of the \texttt{CREATE VIEW} and 93.76\% of the \texttt{INSERT} statements. 
From this, we could populate 6,407 tables with at least one row, 
with a median of 5 rows and a mean of 100.59 rows per table.

To analyze string usage in the examined repositories, we executed the scraped \texttt{SELECT} statements.
Since queries are rarely written as static strings in the code, we first had to preprocess them by 
replacing variables and applying minor transformations to make them executable in DuckDB. This way, we where 
able to execute 27.59\% of the \texttt{SELECT} statements resulting in a 
total number of 56,222 statements we could retrieve query plans for. 

To then be able to track the usage of columns through the operators, we implemented a special 
explain function\footnote{See \url{https://todo}} that returns not only information on 
the operators used to execute the queries but also the details on the expressions these use. 